\documentclass[11pt,a4paper]{article}
% ============================================================
% Préambule commun -- Module M114
% Mathématiques appliquées à la gestion
% Licence MGE -- Université Mundiapolis Business School
% ============================================================
\usepackage[a4paper,top=2.9cm,bottom=2.5cm,left=2.3cm,right=2.3cm,headheight=1.85cm,headsep=8pt]{geometry}
\usepackage[utf8]{inputenc}
\usepackage[T1]{fontenc}
\usepackage{lmodern}
\usepackage[french]{babel}
\usepackage{amsmath,amssymb,mathtools,amsthm}
\usepackage{graphicx}
\usepackage[export]{adjustbox}
\usepackage[dvipsnames]{xcolor}
\usepackage{titlesec}
\usepackage{enumitem}
\usepackage{fancyhdr}
\usepackage{tcolorbox}
\tcbuselibrary{skins,breakable}
\usepackage{array}
\usepackage{booktabs}
\usepackage{tabularx}
\usepackage{multirow}
\usepackage{tikz}
\usetikzlibrary{arrows.meta,calc,positioning}
\usepackage{csquotes}
\usepackage[hidelinks]{hyperref}

% ---------- Couleurs Mundiapolis ----------
\definecolor{mundiblue}{HTML}{0070C0}
\definecolor{mundidark}{HTML}{123B5E}
\definecolor{mundigray}{HTML}{5A5A5A}
\definecolor{munditeal}{HTML}{1B6B6B}
\definecolor{mundired}{HTML}{B03A2E}

% ---------- Titres de sections ----------
\titleformat{\section}
  {\normalfont\Large\bfseries\color{mundidark}}
  {\thesection}{0.6em}{}
\titleformat{\subsection}
  {\normalfont\large\bfseries\color{mundiblue}}
  {\thesubsection}{0.6em}{}
\titleformat{\subsubsection}
  {\normalfont\normalsize\bfseries\color{mundigray}}
  {\thesubsubsection}{0.6em}{}

% ---------- Logos Honoris (gauche) / Mundiapolis (droite) ----------
\graphicspath{{logos/}{./}}
\newcommand{\logoheight}{1.35cm}
% Même hauteur + valign=c : les deux logos sont calés sur le même axe horizontal.
% On ne fixe PAS la largeur : les formats d'origine (Honoris / Mundiapolis) sont différents.
\newcommand{\putlogo}[1]{%
  \includegraphics[height=\logoheight,keepaspectratio,valign=b]{#1}%
}
\newcommand{\headerlogos}{%
  \noindent\makebox[\textwidth][s]{%
    \putlogo{honoris.jpg}\hfill
    \putlogo{munidia.png}%
  }%
}

% ---------- En-tête / pied de page ----------
\pagestyle{fancy}
\fancyhf{}
\renewcommand{\headrulewidth}{0pt}
\renewcommand{\footrulewidth}{0.4pt}
\fancyhead[L]{\putlogo{honoris.jpg}}
\fancyhead[R]{\putlogo{munidia.png}}
\fancyfoot[L]{\small\color{mundigray} Mathématiques appliquées à la gestion -- S1 MGE}
\fancyfoot[C]{\small\color{mundigray}\thepage}
\fancyfoot[R]{\small\color{mundigray} Pr.\ Rajae Malek}
\fancypagestyle{plain}{%
  \fancyhf{}%
  \renewcommand{\headrulewidth}{0pt}%
  \renewcommand{\footrulewidth}{0.4pt}%
  \fancyhead[L]{\putlogo{honoris.jpg}}%
  \fancyhead[R]{\putlogo{munidia.png}}%
  \fancyfoot[C]{\small\color{mundigray}\thepage}%
}

% ---------- Boîtes pédagogiques ----------
\newtcolorbox{definitionbox}[1][]{
  breakable, enhanced, colback=blue!4, colframe=mundiblue,
  boxrule=0.6pt, arc=1mm, left=6pt, right=6pt, top=4pt, bottom=4pt,
  fonttitle=\bfseries, coltitle=white, colbacktitle=mundiblue,
  title={Définition #1}
}
\newtcolorbox{propbox}[1][]{
  breakable, enhanced, colback=green!5, colframe=OliveGreen!70!black,
  boxrule=0.6pt, arc=1mm, left=6pt, right=6pt, top=4pt, bottom=4pt,
  fonttitle=\bfseries, coltitle=white, colbacktitle=OliveGreen!70!black,
  title={Propriété #1}
}
\newtcolorbox{theoremebox}[1][]{
  breakable, enhanced, colback=blue!6, colframe=mundidark,
  boxrule=0.8pt, arc=1mm, left=6pt, right=6pt, top=4pt, bottom=4pt,
  fonttitle=\bfseries, coltitle=white, colbacktitle=mundidark,
  title={Théorème #1}
}
\newtcolorbox{exemplebox}[1][]{
  breakable, enhanced, colback=yellow!8, colframe=orange!80!black,
  boxrule=0.6pt, arc=1mm, left=6pt, right=6pt, top=4pt, bottom=4pt,
  fonttitle=\bfseries, coltitle=black, colbacktitle=orange!25,
  title={Exemple #1}
}
\newtcolorbox{gestionbox}[1][Application en gestion]{
  breakable, enhanced, colback=orange!6, colframe=orange!80!black,
  boxrule=0.6pt, arc=1mm, left=6pt, right=6pt, top=4pt, bottom=4pt,
  fonttitle=\bfseries, coltitle=black, colbacktitle=orange!30,
  title=#1
}
\newtcolorbox{exobox}[1][]{
  breakable, enhanced, colback=gray!5, colframe=mundidark,
  boxrule=0.6pt, arc=1mm, left=6pt, right=6pt, top=4pt, bottom=4pt,
  fonttitle=\bfseries, coltitle=white, colbacktitle=mundidark,
  title={Exercice #1}
}
\newtcolorbox{corrigebox}[1][]{
  breakable, enhanced, colback=gray!3, colframe=gray!60!black,
  boxrule=0.6pt, arc=1mm, left=6pt, right=6pt, top=4pt, bottom=4pt,
  fonttitle=\bfseries, coltitle=white, colbacktitle=gray!60!black,
  title={Corrigé -- #1}
}
\newtcolorbox{methodebox}[1][]{
  breakable, enhanced, colback=teal!4, colframe=munditeal,
  boxrule=0.6pt, arc=1mm, left=6pt, right=6pt, top=4pt, bottom=4pt,
  fonttitle=\bfseries, coltitle=white, colbacktitle=munditeal,
  title={Méthode #1}
}
\newtcolorbox{attentionbox}[1][Piège classique]{
  breakable, enhanced, colback=red!4, colframe=mundired,
  boxrule=0.6pt, arc=1mm, left=6pt, right=6pt, top=4pt, bottom=4pt,
  fonttitle=\bfseries, coltitle=white, colbacktitle=mundired,
  title=#1
}
\newtcolorbox{planbox}[1][Plan de la séance (3 heures)]{
  breakable, enhanced, colback=mundidark!4, colframe=mundidark,
  boxrule=0.6pt, arc=1mm, left=6pt, right=6pt, top=4pt, bottom=4pt,
  fonttitle=\bfseries, coltitle=white, colbacktitle=mundidark,
  title=#1
}
\newtcolorbox{objectifbox}[1][Objectifs]{
  breakable, enhanced, colback=mundiblue!4, colframe=mundiblue,
  boxrule=0.6pt, arc=1mm, left=6pt, right=6pt, top=4pt, bottom=4pt,
  fonttitle=\bfseries, coltitle=white, colbacktitle=mundiblue,
  title=#1
}
\newtcolorbox{remarquebox}[1][Remarque]{
  breakable, enhanced, colback=gray!6, colframe=mundigray,
  boxrule=0.5pt, arc=1mm, left=6pt, right=6pt, top=4pt, bottom=4pt,
  fonttitle=\bfseries, coltitle=white, colbacktitle=mundigray,
  title=#1
}

% ---------- Macros ----------
\newcommand{\N}{\mathbb{N}}
\newcommand{\R}{\mathbb{R}}
\newcommand{\un}{(u_n)}
\newcommand{\vn}{(v_n)}
\newcommand{\wn}{(w_n)}
\newcommand{\limn}{\lim_{n\to +\infty}}
\newcommand{\MAD}{\,\mathrm{MAD}}
\newcommand{\eps}{\varepsilon}

\DeclareMathOperator{\card}{card}

\setlength{\parindent}{0pt}
\setlength{\parskip}{4pt}
\renewcommand{\arraystretch}{1.25}
\setlist[enumerate]{leftmargin=1.4cm, itemsep=3pt}
\setlist[itemize]{leftmargin=1.2cm, itemsep=2pt}

% Page de garde générique
% \pagedegarde{sous-titre}{nature du document}
\newcommand{\pagedegarde}[2]{%
\begin{titlepage}
\hypersetup{pageanchor=false}
\headerlogos\\[0.7cm]
\centering
{\Huge\bfseries\color{mundiblue} UNIVERSITÉ MUNDIAPOLIS}\\[0.45cm]
{\color{mundigray}\large Business School}\\[0.25cm]
{\color{mundigray}\rule{0.5\textwidth}{0.4pt}}\\[1.1cm]
{\Large\color{mundidark} #2}\\[0.45cm]
{\Huge\bfseries\color{mundiblue} Mathématiques appliquées\\[0.15cm] à la gestion}\\[0.7cm]
{\Large\itshape\color{mundidark} #1}\\[1.4cm]
\begin{tabular}{@{}r l@{}}
\textbf{Module :} & M114 \\[4pt]
\textbf{Filière :} & \begin{tabular}[t]{@{}l@{}}Licence en Management\\ et Gestion des Entreprises (MGE)\end{tabular} \\[10pt]
\textbf{Niveau -- Semestre :} & S1 -- 1\up{ère} année \\[4pt]
\textbf{Enseignant(e) :} & Pr.\ Rajae Malek \\[4pt]
\textbf{Année universitaire :} & 2026 -- 2027 \\
\end{tabular}
\vfill
{\color{mundigray}\rule{0.7\textwidth}{0.4pt}}\\[0.3cm]
{\small\color{mundigray} Université Mundiapolis -- Tous droits réservés}
\end{titlepage}
\hypersetup{pageanchor=true}
}


\begin{document}

\pagedegarde{Chapitre 1 -- Les suites réelles}{Feuille d'exercices d'entraînement}

\section*{Présentation}

Cette feuille complète le TD de 3 heures. Elle est destinée au travail personnel (révisions, préparation du contrôle). Les exercices sont classés par thème ; un exercice de synthèse clôt la feuille. Les corrigés détaillés figurent dans \texttt{Chapitre1\_Suites\_reelles\_Corrections.pdf}.

\begin{objectifbox}[Indications de travail]
\begin{itemize}
  \item Exercices 1 à 4 : maîtriser le vocabulaire et les théorèmes de convergence.
  \item Exercices 5 à 7 : suites particulières et finance (cœur du module MGE).
  \item Exercice 8 : récurrence d'ordre 2.
  \item Exercice 9 : problème de synthèse type contrôle.
\end{itemize}
Une copie rédigée de l'exercice 9 constitue un excellent entraînement à l'examen.
\end{objectifbox}

% ------------------------------------------------------------
\section*{Exercice 1 -- Vocabulaire et premières études}

Pour chacune des suites suivantes, préciser si elle est majorée, minorée, bornée, monotone, convergente. Calculer la limite le cas échéant.

\begin{enumerate}
  \item $a_n = \dfrac{n}{n+1}$
  \item $b_n = n^2 - 5n$
  \item $c_n = \left(\dfrac{1}{2}\right)^n$
  \item $d_n = (-1)^n + \dfrac{1}{n+1}$
  \item $e_n = \dfrac{2n-1}{n+3}$ \quad (exercice 11.1 du cours)
\end{enumerate}

% ------------------------------------------------------------
\section*{Exercice 2 -- Limites et formes indéterminées}

Déterminer les limites suivantes.

\begin{enumerate}
  \item $\displaystyle\limn \dfrac{n^2 - 3n + 1}{2n^2 + n}$
  \item $\displaystyle\limn \dfrac{5n+3}{n^2 + 1}$
  \item $\displaystyle\limn \bigl(\sqrt{n+1}-\sqrt{n}\bigr)$
  \item $\displaystyle\limn n\bigl(\sqrt{1+\dfrac{1}{n}} - 1\bigr)$
  \item $\displaystyle\limn \dfrac{4^n - 2^n}{4^n + 3^n}$
\end{enumerate}

% ------------------------------------------------------------
\section*{Exercice 3 -- Encadrements et suites extraites}

\begin{enumerate}
  \item Soit $u_n = \dfrac{n + (-1)^n}{n+2}$. Encadrer $u_n$ et en déduire sa limite.
  \item Soit $v_n = \cos(n\pi) + \dfrac{1}{2^n}$. Étudier $(v_{2n})$ et $(v_{2n+1})$. La suite $\vn$ converge-t-elle ?
  \item On suppose $\un$ convergente vers $\ell$. Montrer que la suite des moyennes de Cesàro $c_n = \dfrac{u_0+\cdots+u_n}{n+1}$ converge aussi vers $\ell$ dans le cas particulier $u_n = \ell + \dfrac{(-1)^n}{n+1}$. (On pourra utiliser le théorème des gendarmes.)
\end{enumerate}

% ------------------------------------------------------------
\section*{Exercice 4 -- Suites adjacentes et récurrence}

\subsection*{A. Suites adjacentes}

On pose $u_n = 2 - \dfrac{1}{2^n}$ et $v_n = 2 + \dfrac{1}{2^n}$.
\begin{enumerate}
  \item Montrer que $\un$ et $\vn$ sont adjacentes.
  \item Déterminer leur limite commune $\ell$.
  \item À partir de quel rang a-t-on $v_n - u_n \leq 10^{-3}$ ?
\end{enumerate}

\subsection*{B. Suite récurrente (exercice 11.4 du cours)}

Soit $u_0 = 0$ et $u_{n+1} = \sqrt{2u_n + 3}$.
\begin{enumerate}
  \item Montrer que pour tout $n$, $0 \leq u_n \leq 3$.
  \item Étudier le sens de variation de $\un$.
  \item En déduire que $\un$ converge et calculer sa limite.
\end{enumerate}

\subsection*{C. Une autre récurrence}

Soit $u_0 = 1$ et $u_{n+1} = \dfrac{u_n + 6}{2}$. Montrer que $\un$ est croissante et majorée par $6$, puis déterminer sa limite.

% ------------------------------------------------------------
\section*{Exercice 5 -- Suites arithmétiques et géométriques}

\begin{enumerate}
  \item Une entreprise constitue une provision : elle verse $12\,000\MAD$ la première année, puis $1\,500\MAD$ de plus chaque année. Combien aura-t-elle versé au total après $8$ ans ?
  \item Un stock de $5\,000$ articles diminue de $12\%$ par an. Exprimer le stock $S_n$ après $n$ années. Au bout de combien d'années le stock sera-t-il inférieur à $1\,000$ articles ?
  \item Un actif de $40\,000\MAD$ se déprécie de $15\%$ par an. Quelle est sa valeur après $6$ ans ? Cette dépréciation est-elle plus rapide que l'amortissement linéaire de $6\,000\MAD$ par an ?
\end{enumerate}

% ------------------------------------------------------------
\section*{Exercice 6 -- Emprunt et tableau d'amortissement}

Une PME emprunte $150\,000\MAD$ au taux annuel $t=5\%$, remboursé par une annuité constante $A = 19\,500\MAD$. On note $D_n$ le capital restant dû après $n$ versements ($D_0 = 150\,000$).

\begin{enumerate}
  \item Écrire la relation de récurrence vérifiée par $(D_n)$ et en donner le terme général.
  \item Compléter le tableau d'amortissement des \textbf{trois premières} années, selon le modèle :

  \medskip
  {\footnotesize
  \noindent
  \begin{tabularx}{\linewidth}{c|>{\centering\arraybackslash}X>{\centering\arraybackslash}X>{\centering\arraybackslash}X>{\centering\arraybackslash}X}
  \toprule
  $n$ & Capital début & Intérêts $t D_{n-1}$ & Amort.\ du capital & Capital fin $D_n$ \\
  \midrule
  1 & $150\,000$ &  &  &  \\
  2 &  &  &  &  \\
  3 &  &  &  &  \\
  \bottomrule
  \end{tabularx}
  }

  \item Vérifier la cohérence avec la formule explicite de $D_3$.
  \item Estimer la durée de l'emprunt (plus petit $n$ tel que $D_n\leq 0$).
\end{enumerate}

% ------------------------------------------------------------
\section*{Exercice 7 -- Placement avec versements périodiques}

Un épargnant place $8\,000\MAD$ le 1\up{er} janvier de l'année 0, puis verse chaque 1\up{er} janvier suivant une somme constante $V = 2\,000\MAD$. Le taux annuel est $t=3\%$. On note $E_n$ l'encours au 1\up{er} janvier de l'année $n$, \textbf{juste après} le versement de l'année $n$.

\begin{enumerate}
  \item Justifier que $E_0 = 8\,000$ et $E_{n+1} = 1{,}03\, E_n + 2\,000$.
  \item Exprimer $E_n$ en fonction de $n$.
  \item Calculer $E_{10}$.
  \item Combien l'épargnant a-t-il versé au total en $10$ ans (versements seuls) ? Quelle est la part des intérêts dans $E_{10}$ ?
\end{enumerate}

% ------------------------------------------------------------
\section*{Exercice 8 -- Récurrence d'ordre 2}

Un fonds obligataire verse un coupon $C_n$ (en MAD) vérifiant
\[
C_{n+2} = C_{n+1} + C_n,
\]
avec $C_0 = 100$ et $C_1 = 150$ (exercice 11.5 du cours).

\begin{enumerate}
  \item Écrire l'équation caractéristique et déterminer ses racines $\varphi$ et $\psi$.
  \item Déterminer $\alpha$ et $\beta$ tels que $C_n = \alpha \varphi^n + \beta \psi^n$.
  \item Calculer $C_2$, $C_3$ et $C_4$ de deux façons (récurrence directe, puis formule).
  \item Justifier que $C_n \to +\infty$ et commenter le réalisme du modèle à long terme.
\end{enumerate}

On donne $\varphi = \dfrac{1+\sqrt{5}}{2} \approx 1{,}618$ (nombre d'or) et $\psi = \dfrac{1-\sqrt{5}}{2} \approx -0{,}618$.

% ------------------------------------------------------------
\section*{Exercice 9 -- Problème de synthèse (type contrôle)}

Une société de location de matériel gère un parc de camionnettes.

\begin{enumerate}
  \item \textbf{Investissement initial.} Elle achète $20$ véhicules à $180\,000\MAD$ l'unité. L'amortissement linéaire retenu est de $22\,500\MAD$ par véhicule et par an. Exprimer la valeur comptable $V_n$ d'un véhicule après $n$ années. Au bout de combien d'années un véhicule est-il totalement amorti ?
  \item \textbf{Recettes.} Le chiffre d'affaires annuel $R_n$ (en milliers de MAD) lié à la location vérifie
  \[
  R_{n+1} = 0{,}85\, R_n + 40, \qquad R_0 = 200.
  \]
  \begin{enumerate}
    \item Interpréter les coefficients $0{,}85$ et $40$.
    \item Exprimer $R_n$ en fonction de $n$ et déterminer $\limn R_n$. Interpréter.
  \end{enumerate}
  \item \textbf{Trésorerie.} On approxime la trésorerie de fin d'année par $T_n = 50 + 2n + \dfrac{(-1)^n}{n+1}$ (en milliers de MAD).
  \begin{enumerate}
    \item Montrer que $T_n \to +\infty$.
    \item La suite $(T_n)$ est-elle monotone ? (on pourra examiner $T_{n+1}-T_n$ ou des extraits).
  \end{enumerate}
  \item \textbf{Demande de location.} Un consultant propose le modèle
  \[
  D_{n+2} = 0{,}5\, D_{n+1} + 0{,}5\, D_n, \qquad D_0 = 80,\ D_1 = 90
  \]
  (nombre de jours de location, en centaines). Déterminer $D_n$ et sa limite.
\end{enumerate}

\vspace{0.6cm}
\begin{center}
{\small\color{mundigray} Fin de la feuille d'exercices}
\end{center}

\end{document}
